\documentclass[12pt,aspectratio=169]{beamer}

\usetheme{tdc}

\title{TDCI data standards}
\subtitle{Status, roadmap, and how to get involved}

\date{
  \texorpdfstring{Joint meeting of \\
  International Transport Energy Modeling (iTEM)\\
  \& Transport Data Commons Initiative (TDCI) \\
  Wed, 18 September 2024}%
  {2024-09-18}}

\author{\texorpdfstring{Paul Natsuo Kishimoto\scriptsize\newline
  \href{mailto:mail@paul.kishimoto.name}%
       {\ttfamily <mail@paul.kishimoto.name>}}%
  {Paul Natsuo Kishimoto <mail@paul.kishimoto.name>}}

\usepackage[
  maxnames = 1,
  style = authoryear,
  giveninits,
  terseinits,
  maxcitenames = 3,
  ]{biblatex}
\addbibresource{all.bib}

\usepackage[table]{xcolor}

\newcommand\checkcell{\cellcolor{tdcyellow0} \checkmark}

\begin{document}

\maketitle

\begin{frame}
\frametitle{Outline}

\tableofcontents

\end{frame}

\section{Why do we need standardization?}

\begin{frame}
\frametitle{Why do we need standardization?}
\framesubtitle{Data from large institutions}

The following slides show Excel-based data file formats from 3 large organizations.

\medskip
Note \emph{where} and \emph{how} each presents (or \emph{does not} present)…
\begin{enumerate}
  \item The \structure{measure} and \structure{units of measure} —i.e. what the numbers represent.
  \item The number, order, and labels of \structure{dimensions} of the data.
  \item The \structure{key values} or labels \emph{along} each of those dimensions.
  \item The \structure{spatial} dimension or scope, including labels for spatial units.
  \item Missing values.
  \item Information about where the data is from, who produced it, when, etc.
\end{enumerate}

\end{frame}

\begin{frame}[plain]
\hspace*{-10mm}
\includegraphics[width=\paperwidth]{2024-09-16_screenshot_jrc-idees-2015.png}
\end{frame}
\begin{frame}[plain]
\hspace*{-10mm}
\includegraphics[width=\paperwidth]{2024-09-16_screenshot_iea-eei.png}
\end{frame}
\begin{frame}[plain]
\hspace*{-10mm}
\includegraphics[width=\paperwidth]{2024-09-16_screenshot_adb-ato.png}
\end{frame}

\begin{frame}
\frametitle{Why do we need standardization?}
\framesubtitle{Data from large institutions}

Each of (1)–(6) is represented \structure{differently} in these 3 file formats!

\bigskip
To be clear:
\begin{itemize}
  \item People worked hard on these data files.
  \item The representations are the result of deliberate decisions.
  Effort was made to put the data in these formats.
  \item The values may have been carefully checked or adjusted for quality.
\end{itemize}

\bigskip
…yet these data are still \structure{not interoperable}.
Where parts of (1)–(6) are \emph{not shown at all}, they are \structure{not reusable} without additional effort to reach a complete description of the data.
\end{frame}

\begin{frame}
\frametitle{Why do we need standardization?}
\framesubtitle{Data from researchers' work}

Researchers also have a bad habit of \structure{inventing new data formats}.
\begin{itemize}
  \item Input, output and intermediate data is often transformed with ‘scripts’.
  \item Output from the first version that ‘works’ is used; attention moves on to other model-building tasks.
\end{itemize}

\medskip
Even when researchers later follow best practice by making data used in/generated from research Findable and Accessible, it is usually not Interoperable or Reusable (“FA” but not “FAIR”).
\begin{itemize}
  \item This increases the time cost if/when the same or other researchers want to, potentially, re-use the data.
\end{itemize}

\end{frame}

\begin{frame}
\frametitle{Why do we need standardization?}
\framesubtitle{Data for and from G-/NTEMs%
\footnote{Global- or national transport-energy models.}}

G-/NTEMs are \structure{data hungry}: they need values for many, varied, multi-dimensional parameters.
At the same time, they are \structure{under-} \structure{determined}: there is less data accessible than parameter values to be set.

\bigskip
Handling input and output data takes up an \structure{unjustifiable share of time} for researchers who build and apply G-/NTEMs.
\begin{itemize}
  \item It \emph{is} good to carefully consider the \emph{validity} and \emph{meaning/ implications} of data.
  \item It \emph{is not} a good use of time to hunt for, struggle to understand, and discover \& fix issues in data—especially if this repeats others' work.
\end{itemize}

\end{frame}

\begin{frame}
\frametitle{Why do we need standardization?}

This is \structure{not} an inevitable, necessary, or justifiable state of practice.\\
We can and must do better.

\medskip
By learning to use data standards, we will save time in order to:
\begin{itemize}
  \item Produce higher-quality models \& research: improve methods, refine scenarios, compare/analyse results more precisely.
  \item Allow more people to join in and advance the work.
  \item Communicate research results to stakeholders and help them understand policy implications.
\end{itemize}

\medskip
These benefits of standardization are visible when looking at research disciplines (e.g. Earth sciences, genetics) that have recognized its importance.
\end{frame}

\section{What do the TDC standards comprise?}
\begin{frame}
\frametitle{What do the TDC standards comprise?}

The standards can be boiled down to one instruction:

\medskip
\begin{quote}
  \usebeamercolor{structure} \large Express your data and metadata using SDMX.
\end{quote}

\bigskip
Everything else is details, and can be found at \url{https://docs.transport-data.org/en/latest/standards.html}

\end{frame}

\subsection{Statistical Data and Metadata eXchange (SDMX)}

\begin{frame}
\frametitle{Statistical Data and Metadata eXchange (SDMX)}
\framesubtitle{What is it?}

\begin{itemize}
  \item An ISO standard (ISO 17369:2013) actively developed since 2005 by a group including the World Bank, IMF, UN, Eurostat, European Central Bank, and others.
  \item Adopted by 30+ national statistical agencies and many more organizations.
  \item Per the name:
    \begin{itemize}
      \item Focused on \structure{exchange} (=interoperability and reusability) via data file formats (XML, JSON, CSV) and web APIs.
      \item Inclusive both of \structure{data} (actual values) and \structure{metadata} (information \emph{about} data: its structure \emph{and} provenance).
    \end{itemize}
  \item \structure{Generalized} and \structure{non-domain-specific} via an \structure{information model} that is inclusive of many kinds of (meta)data.
\end{itemize}
\end{frame}

\begin{frame}[fragile,allowframebreaks]
\frametitle{Example usage: IAEG-SDGs}
\framesubtitle{Inter-agency and Expert Group on Sustainable Development Goal (SDG) Indicators}

This group produced metadata to be used with SDG-related data:

→ \href{https://unstats.un.org/sdgs/iaeg-sdgs/sdmx-working-group/}{unstats.un.org/sdgs/iaeg-sdgs/sdmx-working-group}

\bigskip
Based on their work one can say, for instance:

\medskip
“\emph{My} data has a \structure{dimension} \verb@SERIES@ that is enumerated by the \structure{code list} \verb@IAEG-SDGs:CL_SERIES(1.18)@.”
\begin{itemize}
  \item This unambiguously identifies the maintainer, ID, and version.
  \item Even if new versions of this code list are released, we know which codes are used in the described data.
\end{itemize}

\framebreak
\medskip
“The key for a value ‘12.3’ includes (…, \verb@SERIES=SI_COV_BENEFITS@, …).”
\begin{itemize}
  \item This ID is short and easily manipulated by data-handling code.
  \item Because the data has a defined structure + is associated with a code list, we can retrieve detailed, structured metadata to support proper use of the values:
  \begin{itemize}
    \item Description: Proportion of population covered by at least one social protection benefit [1.3.1]
    \item IndicatorTitle: Proportion of population covered by social protection floors/systems, by sex, distinguishing children, unemployed persons, older persons, persons with disabilities, pregnant women…
  \end{itemize}
\end{itemize}
…the same applies to all other concepts and dimensions in SDG-related data.
\end{frame}

\begin{frame}
\frametitle{Example usage: Eurostat Data Browser}
\framesubtitle{Presentation adapts to \textbf{any} structure, via SDMX metadata}
\includegraphics[width=\textwidth]{2024-09-16_screenshot_eurostat.png}
\end{frame}

\subsection{“TDC standards” = a common way of using SDMX}
\begin{frame}
\frametitle{“TDC standards” = a common way of using SDMX}

The labels “public data”, “community data”, “TDC formatted/ compatible”, and “TDC Harmonized” directly give the degree to which (meta)data have been made accessible and interoperable \& quality standards applied.

\medskip
\begin{tabular}{l|c|c|c|c}
Characteristic of (meta)data
  & \cellcolor{tdcgreen0}\includegraphics{tdc/globe.pdf}
  & \cellcolor{tdcgreen0}\includegraphics{tdc/database.pdf}
  & \cellcolor{tdcgreen0}\includegraphics{tdc/check-circle.pdf}
  & \cellcolor{tdcgreen0}\includegraphics{tdc/check-shield.pdf} \\
\hline
Metadata exist                            & ?? & \checkcell & \checkcell & \checkcell \\
Metadata for TDC attributes (in SDMX)     & —  & \checkcell & \checkcell & \checkcell \\
Data in SDMX formats (possibly others)    & —  & —          & \checkcell & \checkcell \\
Uses shared concepts, structures, and IDs & —  & —          & —          & \checkcell \\
TDC quality checks \& adjustments applied & —  & —          & —          & \checkcell \\
Freely accessible                         & ?? & \checkcell & ??         & \checkcell \\
\hline
\end{tabular}

\end{frame}

\begin{frame}[fragile,allowframebreaks]
\frametitle{Clarifying points}

\structure{TDCI only does coordination.}
\begin{itemize}
  \item Initial standards—like “Such-and-such concept/dimension \textbf{should} have the ID ‘\verb@VEHICLE_TYPE@’,”—are merely an \structure{observation of common practice}.
  \item TDCI convenes members/stakeholders so \emph{they} can discuss the most useful sets of labels, metadata attributes, etc.
  \item iTEM participants are a major, important subset of this community.
  \item After the community decides, TDCI helps with:
    \begin{itemize}
      \item Tools that help apply/use the agreed standards.
      \item Documentation, resources, and a focal point for more users to discover them (and potentially also join the community to contribute).
    \end{itemize}
\end{itemize}

\framebreak
\structure{No one is obliged to use, or do, anything.}
\begin{itemize}
  \item TDC will handle metadata for “public data” and “community data”, and serve as a repository for community data \emph{files} —these will simply be handled as black boxes/blobs.
  \item “TDC compatible” (meta)data can be shared without using shared concepts and code lists.

  \item Often it is \emph{better} to \structure{describe original data \emph{as they are}}.

  If Data Set A uses a label “mutatu”, and Data Set B “tuk tuk”, this conveys what people think is important, accurate, and meaningful. This is more informative than an outsider's choice to lump these under e.g. “other”.
  \item Researchers may use new labels and categories that in developing \structure{new methods and empirical findings}—also good. TDC metadata can express how these relate to existing codes.
\end{itemize}

\framebreak
\structure{No one is prohibited from doing anything.}
\begin{itemize}
  \item If researchers/modelers/others have tools that work with idiosyncratic input/output formats—they can continue to use those.
  \item Models/tools can I/O SDMX directly \emph{or} SDMX can be converted to/from specific other formats.
  \item These data handling workflows will be \emph{easier to develop} and maintain, and \emph{more likely to be reusable} because they will be N:1 / 1:N, rather than many-to-many.
\end{itemize}

\end{frame}

\section{Standardization is a process}

\begin{frame}
\frametitle{Standardization is a process}

The process includes:
\begin{description}
  \item [Learning] what the standards are.
  \item [Applying] them to an increasing degree and extent.
  \item [Gaining skill] and facility in applying standards—including through building assistive tools.
  \item [Developing] standards to ensure they continue to provide benefits.
\end{description}

\bigskip
\begin{itemize}
  \item Each individual/team can \structure{take steps incrementally}, as their resources and interests allow.
  \item As they do, they will experience \structure{progressively greater benefit} (time savings, etc.) \emph{and} provide an improved resource to the community (network effects).
\end{itemize}

\end{frame}

\subsection{How to get involved—roles and activities}

\begin{frame}
\frametitle{How to get involved with standardization}

Both iTEM and TDCI \structure{as organizations/groups} and their \structure{individual members/participants} have many options:

\smallskip
\begin{tabular}{l|cc|cc}
 & \multicolumn{2}{c|}{TDCI} & \multicolumn{2}{c}{iTEM} \\
Activity & Mem & Org & Mem & Org \\
\hline
Use TDC-formatted (meta)data      & \checkcell &            & \checkcell &            \\
Produce TDC-formatted (meta)data  & \checkcell & \checkcell & \checkcell & \checkcell \\
Improve tools for the above       & \checkcell & \checkcell & \checkcell & \checkcell \\
\hline
Propose additions/improvements    & \checkcell &            & \checkcell &            \\
Convene discussions of standards  &            & \checkcell &            & \checkcell \\
\hline
Set general quality criteria      & \checkcell &            & \checkcell &            \\
Set scientific quality criteria   &            &            & \checkcell &            \\
\hline
Encourage adoption in research    &            &            & \checkcell & \checkcell \\
Encourage adoption by *GOs        &            & \checkcell &            &            \\
\hline
\end{tabular}

\end{frame}

\subsection{Next steps and vision}

\begin{frame}[allowframebreaks]
\frametitle{Next steps and vision}

\structure{2024–2025}
\begin{itemize}
  \item iTEM Open Data is refreshed using TDC tools.
  \item Transparent, repeatable quality checks and fixes developed by iTEM support first examples of TDC Harmonized.
\end{itemize}

\bigskip
\structure{2025} at iTEM8:
\begin{itemize}
  \item iTEM modelers' outputs (a subset) are compared via data and structures shared on TDC.
  \item Some modelers are using TDC-formatted data directly for aligned inputs, or expressing inputs to allow adjusted comparison of outputs.
  \item TDC standards are updated to be inclusive of all input/output concepts, dimensions, measures used by iTEM modelers.
\end{itemize}

\framebreak
\structure{2027–2028} as part of the Working Group III contribution to IPCC AR7:
\begin{itemize}
  \item A richer set of measures (indicators) from a wider range of models is available and comparable via standardized (meta)data.
  \item This process begins early and in public, allowing better science through iteration and broad participation.
  \item The current and future mobility of people in low- and middle-income countries can be assessed using an expanded evidence base, using data coming directly from practice.
\end{itemize}

\end{frame}

\section*{Thank you!}

% \appendix

% \begin{frame}
% \frametitle{References}
% \printbibliography
% \end{frame}

\end{document}
